\chapter[State of the Art]{State of the Art and Technological\\Background}\label{chap:soa}
\begin{spacing}{1.2}
\minitoc
\thispagestyle{MyStyle}
\end{spacing}
\newpage

\section*{INTRODUCTION}
\addcontentsline{toc}{section}{INTRODUCTION}
The system proposed in this report operates at the intersection of three technological domains: industrial quotation and Request for Quotation (RFQ) management; microservices and layered architectural patterns; and conversational agents powered by Large Language Models (LLMs). This chapter critically reviews each landscape to identify the specific operational gap the proposed platform addresses, rather than conducting an exhaustive survey of any single field. Section~\ref{sec:positioning} synthesizes these findings and defines the strategic positioning of the proposed platform.

\section{Industrial RFQ Systems and the Lifecycle Visibility Gap}\label{sec:rfq_systems}
Three categories of commercial software predominantly support industrial quotation. \textbf{Enterprise Resource Planning (ERP)-integrated quoting modules}, exemplified by SAP CPQ~\cite{sap_cpq_product}, embed quote generation within the transactional and master-data backbone of the enterprise. \textbf{Standalone Configure, Price, Quote (CPQ) systems}, exemplified by Salesforce CPQ~\cite{salesforce_cpq_docs}, focus on configuration, pricing, and proposal generation, offering advanced support for configuration logic and discounting rules. \textbf{Custom spreadsheet-based estimation} remains prevalent in domains where the cost model relies heavily on uncodified expert judgment, as exemplified by the GHI estimation workbook discussed in Chapter~\ref{chap:context}. Business process management literature~\cite{dumas2018bpm} frames the RFQ activity as a coordinated orchestration spanning multiple roles and decision points, distinct from any isolated price calculation.

Each category optimizes a specific subset of the quotation activity. None of these paradigms adequately address the specific operational gap identified during the \clientCompany audit: cross-departmental lifecycle coordination, executive-level visibility, capitalization on historical data, early identification of feasibility risks, and the structured capture of interim decisions. The primary operational bottleneck resides not within the estimation workbook, but in its surrounding lifecycle. The proposed platform is designed to provide the combination of capabilities required to address this lifecycle: RFQ orchestration, deterministic document intelligence, evidence-grounded conversational access, and a foundational architecture for future predictive learning.

\section{Microservices Architecture and Project-Specific Layering}\label{sec:microservices_layering}
At the architectural level, the proposed system is structured as five responsibility-segregated microservices~\cite{newman2021microservices}, decomposed along the functional boundaries of the RFQ lifecycle: operational lifecycle management, document and workbook intelligence, conversational assistance, user interaction, and identity and access management. This detailed decomposition is presented in Chapter~\ref{chap:architecture} (Section~\ref{sec:bacab_pattern}). Internally, each backend service implements the \textbf{BACAB pattern}, a project-specific synthesis of the traditional three-tier architecture and the Ports-and-Adapters paradigm~\cite{cockburn2005hexagonal}, informed by clean-architecture layering principles~\cite{martin2017cleanarch}. This pattern separates entry points (routes), business orchestration (controllers), persistence access (datasources), external communication (connectors), and data representation (translators), supported by models, configuration, and utilities layers. This architecture ensures lifecycle traceability, separates operational source-of-truth from derived intelligence, and facilitates independent service evolution; it is conceptualized in Chapter~\ref{chap:architecture} (Section~\ref{sec:bacab_pattern}) and implemented in Chapter~\ref{chap:implementation} (Section~\ref{sec:impl_env}).

\section{Conversational AI Patterns for Evidence-Grounded Systems}\label{sec:conversational_ai_patterns}
The conversational layer of the proposed platform leverages the recent advancements in Large Language Models~\cite{brown2020gpt3}. The four design patterns reviewed below share this foundational technology; however, they differ significantly in the constraints they impose on model behavior. The critical objective is determining the optimal pattern for an industrial context characterized by asymmetric costs of error.

\textbf{Naive Retrieval-Augmented Generation (RAG)}~\cite{lewis2020rag} appends retrieved documents to the prompt; while suitable for question-answering over static corpora, this pattern lacks structural defenses to prevent the model from hallucinating beyond the retrieved context. \textbf{LLM agents with autonomous tool calling}, exemplified by ReAct~\cite{yao2023react}, grant the model autonomy over tool selection and execution sequence. While appropriate for exploratory tasks, this autonomy poses a structural risk in environments where retrieved values inform binding contractual quotations. \textbf{Constrained tool-calling chatbots} leveraging provider function-calling APIs~\cite{anthropic_tool_use_docs} retain execution control within the application code, yet still allow the model to influence policy via intent classification. \textbf{Trust-bounded or policy-as-configuration systems}, the emerging fourth paradigm, delegate all safety-critical decisions (e.g., tool authorization, data access control, refusal triggers, context injection) to deterministic platform components rather than the probabilistic language model. While a canonical designation for this architectural family remains to be standardized, it represents the architectural commitment of the proposed platform, justified by the safety and grounding imperatives developed in Section~\ref{sec:hallucination} and operationalized in Chapter~\ref{chap:architecture} (Section~\ref{sec:trust_boundary}).

\section{Grounding, Hallucination, and Safety Constraints}\label{sec:hallucination}
The deployment of any language-model-based conversational agent necessitates the mitigation of \textbf{hallucination}: the generation of factually incorrect yet locally coherent and stylistically plausible content. The survey by Ji et al.~\cite{ji2023hallucination} establishes hallucination as an intrinsic structural property of the generative paradigm rather than a transient anomaly; consequently, mitigation mechanisms must be engineered into the surrounding system architecture rather than expected to emerge from the model inherently. A corollary constraint is \textbf{context dilution}, wherein relevant evidence is obfuscated by extraneous material within an extensive prompt. This constraint necessitates the implementation of the per-path field whitelist developed in Chapter~\ref{chap:architecture} (Section~\ref{sec:pipeline_vertical_slicing}).

While Retrieval-Augmented Generation mitigates hallucination, it does not eradicate it; the model remains susceptible to drifting beyond the provided evidence. A system relying exclusively on RAG provides insufficient structural defense against such drift. Therefore, the proposed copilot integrates four complementary mechanisms: \textbf{deterministic guardrails} evaluate the generated response prior to delivery (enforcing evidence grounding, scope adherence, and shape conformance); an \textbf{LLM-as-Judge layer}~\cite{zheng2023judge} executes final linguistic verification (detecting fabrication, scope drift, and forbidden inferences) that deterministic rules cannot reliably capture; \textbf{structured outputs}~\cite{anthropic_structured_outputs_docs} constrain the model to generate strictly schema-conformant JSON; and \textbf{policy-\allowbreak as-\allowbreak configuration}~\cite{bai2022constitutional} encodes safety policies as declarative data rather than procedural code, enabling auditable scope modifications without necessitating pipeline redeployments. This combination forms the foundation of the trust-boundary architecture, conceptualized in Chapter~\ref{chap:architecture} (Section~\ref{sec:trust_boundary}) and validated in Chapter~\ref{chap:validation} (Section~\ref{sec:trust_boundary_verification}).

\section{Deterministic Document Intelligence for Engineering Artifacts}\label{sec:deterministic_intelligence}
The estimation workbook and Material Requisition package, introduced in Chapter~\ref{chap:context}, are structured artifacts whose data parameterizes the cost models. Two primary paradigms exist for extracting information from such documents. \textbf{Deterministic parsing} leverages inherent document structure: spreadsheets with predefined layouts are parsed via cell coordinates, and documents with standardized section hierarchies are parsed by identifiers, yielding structured data without generative interpretation. Mature parsing ecosystems exist for standard formats. \textbf{LLM-based document question answering} submits document content to a language model to extract specific values. While highly adaptable to document variability, this approach is inherently vulnerable to hallucination (Section~\ref{sec:hallucination}), a critical flaw when extracting numeric values that cannot be readily verified visually by users.

The architectural selection between these paradigms is dictated by the cost-of-error posture (Chapter~\ref{chap:context}): an erroneous value propagated into a binding quotation incurs significantly higher financial risk than a missing value that safely triggers a manual review. Consequently, the intelligence service adopts a \textbf{deterministic-parsing-first strategy}: parsers extract structured profiles, and values failing extraction thresholds are flagged as anomalies requiring human intervention rather than being probabilistically approximated. This design principle is the central contribution of the document-intelligence layer, demonstrated by the parser implementations detailed in Chapter~\ref{chap:implementation} (Section~\ref{sec:impl_intelligence}).

These artifacts are strictly governed by industrial standards. Pressure vessel construction complies with the ASME Boiler and Pressure Vessel Code, Section VIII Division 1~\cite{asme_bpvc_viii_div1_2025}; shell-and-tube heat exchangers adhere to the standards of the Tubular Exchanger Manufacturers Association~\cite{tema_standards_11ed}. The intelligence service's parsers are engineered to systematically extract standard references within the Material Requisition package, providing a data foundation that supports the later identification of standards-driven cost implications, addressing a critical operational delay identified during the audit (PP-12).

\section{Positioning of the Proposed Platform}\label{sec:positioning}
The preceding sections have reviewed the systemic categories against which the proposed platform is positioned: ERP-integrated quoting modules, standalone CPQ systems, custom spreadsheet-based estimation, the four paradigms of LLM-powered conversational architectures, and document-intelligence methodologies. These categories are not mutually exclusive competitors; rather, each addresses a distinct facet of the industrial quotation process. The core contribution of the proposed platform resides in synthesizing capabilities tailored to the specific operational gap identified during the audit, rather than claiming superiority in domains natively addressed by incumbent systems.

The proposed platform is formally positioned as an \textbf{RFQ Lifecycle Intelligence Platform integrating a trust-bounded conversational layer}. It synthesizes lifecycle awareness derived from Business Process Management (BPM) principles (Section~\ref{sec:rfq_systems})~\cite{dumas2018bpm}; multi-actor coordination structured via Domain-Driven Design~\cite{evans2003ddd} (Chapter~\ref{chap:architecture}, Section~\ref{sec:manager_domain}); robust evidence grounding mechanisms addressing safety constraints~\cite{ji2023hallucination, bai2022constitutional, zheng2023judge}; strict scope-bounded refusal enforced by structured-output discipline~\cite{anthropic_structured_outputs_docs} and deterministic guardrails; policy-as-configuration auditability characteristic of trust-bounded architectures; and a risk-averse, deterministic-parsing-first strategy for engineering documents. Chapter~\ref{chap:architecture} formulates the architecture enabling this synthesis; Chapter~\ref{chap:implementation} details its technical implementation; and Chapter~\ref{chap:validation} validates the architectural invariants through empirical evaluation.

\section*{CONCLUSION}
\addcontentsline{toc}{section}{CONCLUSION}
This chapter has evaluated the three technological landscapes pertinent to the proposed platform, isolating the operational gaps each paradigm leaves unresolved in the context of the audit findings. The platform is strategically positioned within this gap as an RFQ Lifecycle Intelligence Platform featuring a trust-bounded conversational layer; it does not aim to replace ERPs, compete with CPQs, act as a generic chatbot, or function as a fully autonomous BoQ generator. The specific architecture realizing this strategic positioning is detailed in the subsequent chapter.
